\chapter{Lua Preprocessing, Macros, and Multi-Pass Layout}
\label{ch:assembly-preprocessing}

A repeated instruction sequence is easy to generate until it contains labels, configuration choices, and size-sensitive branches.
Copying its text twice can define the same private label twice.
Evaluating a size expression too early can freeze a value that later instruction expansion invalidates.
Allowing arbitrary scripts to query and mutate a half-built symbol table can make output depend on which assembly pass happened to run first.
A useful extensible assembler therefore needs phase boundaries as carefully as it needs an expression parser.

The case study places Lua before assembly.
Lua computes a finite assembly stream; the assembler then interprets that stream under its own symbol, size, and encoding rules.
This is a concrete form of the macro-expansion phase discussed in assembler literature~\cite{salomon-1993-assemblers}.
It also creates an opportunity to distinguish executable configuration from the later fixed-point problem of layout.

\section{Specify what preprocessing can know}
\label{sec:ch55-phases}

Before assembly begins, a Lua script can know its arguments, constant tables, selected feature flags, and the text it has emitted.
It cannot know the final address of a symbol in another object.
In our interface it also cannot query the assembler's current section counter, because the assembler has not parsed the emitted stream yet.
These restrictions follow directly from phase order.
They prevent an expression such as \enquote{emit this instruction only if the later function is nearby} from accidentally becoming an undocumented dependency cycle.

Figure~\ref{fig:ch55-phases} puts the boundaries in order.
Each stage has an input that is complete enough for its decisions.
An implementation can add later callbacks, but then it must define what can change, how dependencies are recorded, and when another iteration is required.
Embedding a general-purpose language does not itself answer those questions.

\begin{figure}[htbp]
\centering
\begin{tikzpicture}[node distance=13mm]
\node[basic block] (lua) {Lua values and functions};
\node[basic block,below=of lua] (text) {completed assembly text};
\node[basic block,below=of text] (frag) {instructions, labels, fixups};
\node[basic block,below=of frag] (obj) {ELF sections and relocations};
\draw[cfg edge] (lua)--node[right,font=\small] {emit}(text);
\draw[cfg edge] (text)--node[right,font=\small] {parse and encode}(frag);
\draw[cfg edge] (frag)--node[right,font=\small] {serialize}(obj);
\end{tikzpicture}
\caption{Preprocessing finishes before assembly state exists, making the available information at each phase explicit.}
\label{fig:ch55-phases}
\end{figure}

Lua functions supply parameterization, loops supply repetition, and tables supply generated data.
There is no need to invent a second integer-expression language for counts already known at preprocessing time.
An address expression remains textual assembly, however: emitting \code{.quad answer+8} preserves the symbol dependency for later phases rather than asking Lua to evaluate it.
The boundary is between known preprocessing values and relocatable program values.

\section{Private macro labels need fresh identities}
\label{sec:ch55-labels}

Suppose a countdown macro emits a label followed by a subtract and conditional branch.
Calling it twice with the same literal label causes a duplicate definition.
Worse, a permissive assembler might bind both branches to the second definition.
The correct expansion requires a fresh private label for every invocation.
This is a simple example of macro hygiene, even though textual assembly macros generally do not provide the full lexical hygiene of language-level syntax systems.

The Lua example maintains a local serial number and emits \code{.loop\_1}, \code{.loop\_2}, and later names.
The prefix is reserved by convention for that macro implementation; a larger generator should encapsulate name creation behind a shared fresh-name service.
Two modules independently resetting the same serial with the same prefix can still collide when their generated text enters one assembly unit.
Figure~\ref{fig:ch55-fresh} contrasts the two-invocation result.

\begin{figure}[htbp]
\centering
\begin{tikzpicture}[node distance=18mm and 22mm]
\node[basic block] (m) {countdown macro};
\node[basic block,below left=of m] (a) {.loop\_1: sub ecx, 1};
\node[basic block,below right=of m] (b) {.loop\_2: sub edx, 1};
\draw[cfg edge] (m)--node[left,font=\small,fill=white,inner sep=2pt] {first call}(a);
\draw[cfg edge] (m)--node[right,font=\small,fill=white,inner sep=2pt] {second call}(b);
\draw[back edge] (a) edge[loop below] node[below,font=\small] {jnz .loop\_1}(a);
\draw[back edge] (b) edge[loop below] node[below,font=\small] {jnz .loop\_2}(b);
\end{tikzpicture}
\caption{Fresh private names keep each generated control-flow cycle inside its own macro invocation.}
\label{fig:ch55-fresh}
\end{figure}

Parameters should also be validated before interpolation.
A register parameter is not an arbitrary assembly fragment if the function promises a register-based countdown.
The case-study assembler ultimately validates the emitted operands; a reusable Lua library can fail earlier with a clearer message about the macro argument.
For raw byte data, emitting one numeric \code{.byte} per source byte avoids quoting ambiguities involving embedded zeros, newlines, quotes, or comment delimiters.
The \path{hello.lua} example uses this strategy for a message written by a Linux system call.

\begin{algorithm}[htbp]
\caption{Expand a counted assembly macro with a private label.}
\label{alg:ch55-macro}
\begin{pseudocode}
\Require valid register $r$, positive count $n$, shared fresh-name state
\Ensure one self-contained countdown sequence with no reused private label
\State $label\gets\Call{FreshName}{countdown}$
\State emit initialization of $r$ to $n$
\State emit the private label definition
\State emit subtraction of one from $r$
\State emit a conditional branch back to $label$ while the result is nonzero
\end{pseudocode}
\end{algorithm}

Algorithm~\ref{alg:ch55-macro} emits a constant number of assembly statements regardless of the runtime count.
An unrolling macro instead emits $O(n)$ statements and changes code size accordingly.
Compile-time repetition and runtime iteration are different transformations, even if both are written using a loop in the generator's implementation.

\section{Cross the Lua/C boundary with explicit representations}
\label{sec:ch55-lua-api}

The assembler registers one C function, \code{emit}, in a Lua state.
It accepts a Lua byte string, appends its bytes to a checked buffer, and adds a newline.
Embedded zero bytes are rejected because the downstream assembly parser uses nul-terminated line strings.
The data-generation interface still represents zero data bytes by emitting the textual directive \code{.byte 0}.
Do not confuse zeros in the source-language transport with zeros requested in the target program.

Lua compilation and execution occur through a protected call.
If the script raises an error, the host prints the Lua error and never publishes an object file.
Lua error values need not be strings, so the error path must not blindly pass a null conversion result to a C string formatter.
The implementation emits a fallback diagnostic for a non-string error value.
These behaviors follow the Lua C API's error and string conventions~\cite{ierusalimschy-2026-lua}.

Table~\ref{tab:ch55-boundaries} lists the decisions at this interface.
The scripts run with Lua's standard libraries and therefore have the same build-time authority as other build scripts.
The tools are not script sandboxes.
The generated-text buffer is bounded, but an arbitrary Lua program can still run forever before emitting anything; a source-size bound is not an execution-time bound.

\begin{table}[htbp]
\caption{The preprocessing boundary distinguishes source transport, target data, and script execution.}
\label{tab:ch55-boundaries}
\centering
\begin{tabular}{p{0.30\linewidth}p{0.57\linewidth}}
\toprule
Value or event & Case-study behavior \\
\midrule
Lua string passed to emit & Copied by explicit byte length; newline appended. \\
NUL inside source text & Rejected before line parsing. \\
Target byte value zero & Expressed as an ordinary numeric data directive. \\
Lua runtime error & Reported before object publication. \\
Assembly encoding error & Reports generated assembly line number. \\
Final linked address & Unavailable to preprocessing; represented symbolically. \\
\bottomrule
\end{tabular}
\end{table}

\section{A listing should explain expansion provenance}
\label{sec:ch55-provenance}

When a generated instruction fails, the programmer needs to know both what the assembler saw and why it was generated.
A complete provenance record might include the emitted line, Lua file and line, macro name, invocation number, and caller chain.
The generated line remains useful because it identifies the exact syntax that the assembler rejected.
The originating Lua location remains useful because that is the editable source.

The current implementation keeps the source script path and generated assembly line number.
Lua's own syntax and execution diagnostics identify script locations.
It does not claim a complete expansion stack or DWARF line-table producer.
An extension could store a provenance record beside every emitted line, using Lua debug information at the callback boundary and an explicit macro-call identifier for nested libraries.
That record belongs to the intermediate representation, not to an unstructured diagnostic string reconstructed after failure.

\section{Variable-size instructions turn layout into a dependency problem}
\label{sec:ch55-relaxation}

The implemented assembler uses near branches with fixed four-byte displacements.
It therefore does not need branch relaxation to determine label offsets.
To understand a production assembler, now consider adding short two-byte jumps and long five-byte jumps.
A branch's allowed form depends on its distance; the distance depends on the sizes of instructions between source and target; those sizes can include other branches.
One pass that selects every shortest-looking form from an initial layout can invalidate its own decisions.

Use two forward branches $A$ and $B$, initially two bytes each.
After $B$, place $124$ fixed bytes, label $T$, eight more fixed bytes, and label $U$.
$A$ targets $T$ and $B$ targets $U$.
Table~\ref{tab:ch55-relaxation} computes the widening sequence.
Short displacements range from $-128$ to $127$.
Initially $A$ fits, but $B$ does not; widening $B$ then pushes $A$ out of range.

\begin{table}[htbp]
\caption{One branch expansion can invalidate a previously fitting branch.}
\label{tab:ch55-relaxation}
\centering
\begin{tabular}{lllllll}
\toprule
Round & Size A & Size B & T & U & A displacement & B displacement \\
\midrule
Initial & $2$ & $2$ & $128$ & $136$ & $126$ & $132$ \\
Widen B & $2$ & $5$ & $131$ & $139$ & $129$ & $132$ \\
Widen A & $5$ & $5$ & $134$ & $142$ & $129$ & $132$ \\
\bottomrule
\end{tabular}
\end{table}

\begin{figure}[htbp]
\centering
\begin{tikzpicture}[x=13mm,y=10mm]
\node[basic block] (a) at (0,0) {A}; \node[basic block] (b) at (2,0) {B};
\node[basic block] (t) at (5,0) {T}; \node[basic block] (u) at (7,0) {U};
\draw[cfg edge] (a) to[bend left=28] node[above,font=\small] {crosses B's encoding}(t);
\draw[cfg edge] (b) to[bend right=28] node[below,font=\small] {must widen first}(u);
\draw[dashed] (b)--node[above,font=\small] {124 bytes}(t);
\draw[dashed] (t)--node[above,font=\small] {8 bytes}(u);
\end{tikzpicture}
\caption{The interval covered by A includes B, so changing B's size changes A's required displacement.}
\label{fig:ch55-relaxation}
\end{figure}

\begin{algorithm}[htbp]
\caption{Find a valid layout by monotonically widening instruction forms.}
\label{alg:ch55-widen}
\begin{pseudocode}
\Require finite ordered forms per instruction, largest forms checked for representability
\Ensure a self-consistent valid layout or explicit unrepresentable-reference error
\State choose the smallest allowed form for each relaxable instruction
\Repeat
  \State recompute offsets, labels, and alignment padding from current forms
  \State collect instructions whose fields do not fit in the current layout
  \If{any failing instruction has no wider supported form}
    \State \Return unrepresentable-reference error
  \EndIf
  \State widen every collected instruction by at least one form
\Until{no instruction was widened}
\State revalidate every final field, then emit bytes and relocations
\end{pseudocode}
\end{algorithm}

Let $F$ be the number of fragments and $W$ the total number of possible form promotions across them.
Algorithm~\ref{alg:ch55-widen} performs at most $W$ changing rounds, plus a final stable round.
With $O(F)$ layout and checking per round, its bound is $O(F(W+1))$.
The proof of termination uses the monotone form indices, not monotone label distances: alignment can cause padding to increase or decrease as preceding instructions grow.
The final validation establishes representability under the final layout.
This algorithm need not produce the globally smallest encoding.
It never revisits a widening choice even if a later padding change would permit a smaller form.

\section{Alignment, literals, and external targets add constraints}
\label{sec:ch55-layout-constraints}

Alignment fragments are functions of their incoming offset, not constants stored forever after their first visit.
For power-of-two alignment $a$, padding is $(-x)\bmod a$, with a result in $[0,a)$.
Increasing an earlier instruction by one byte can reduce that padding by one and leave later addresses unchanged.
This is why a dependency graph over fragments is more informative than treating every byte increase as a uniform shift of all following labels.

Literal pools introduce another kind of reachability constraint.
An instruction that loads a constant from a nearby pool needs the pool within its displacement range.
Pool placement can require an inserted branch around embedded data and can change the distances of other branches.
Such a pass must preserve fall-through behavior as well as field range.
The x86-64 case study avoids literal pools, but the representation of deferred relationships generalizes to them.

An external target has no known same-object distance.
The assembler can choose a conservative form whose relocation the linker understands, or emit a relaxable sequence with sufficient metadata for link-time rewriting.
It cannot honestly certify a short external branch by assuming that separately compiled functions will happen to be adjacent.
Veneers, thunks, and linker relaxation are target-dependent solutions discussed again in Chapter~\ref{ch:link-layout}.

\section{Reproducibility starts before byte emission}
\label{sec:ch55-reproducibility}

Lua table iteration with \code{pairs} does not define a portable stable order for arbitrary keys.
If emission order matters, use a dense numeric array with an explicit numeric loop or sort a collected key list.
Avoid incorporating time, random values, or ambient paths into bytes unless they are declared build inputs.
Two deterministic assemblers cannot repair nondeterministic source generation that has already changed the instruction order.

The same principle applies to symbol-table ordering and diagnostics.
Stable source order and explicit local/global partitioning make output repeatable without depending on a hash table's bucket order.
Tests compare complete ELF bytes across compiler configurations and repeat generation to catch accidental state leakage.
Determinism is a property of the entire pipeline and its input policy, not only a promise made by the final file writer.

\section{Exercises}
\label{sec:ch55-exercises}
\begin{enumerate}
\item\label{ex:ch55-fresh} Expand two countdown invocations with counts three and five. List their private names and explain which shared state guarantees uniqueness within the assembly unit.
\item\label{ex:ch55-bytes} Give the emitted numeric data directives for the three-byte Lua string containing ASCII A, zero, and newline.
\item\label{ex:ch55-rounds} Recompute Table~\ref{tab:ch55-relaxation} after each widening. Explain why A's displacement stays 129 when A itself grows.
\item\label{ex:ch55-align} Compute padding to alignment 16 at offsets 29, 30, and 32. What happens to the aligned destination after a one-byte growth from 29 to 30?
\item\label{ex:ch55-bound} There are 20 branches with two forms each and 200 fragments. Give the maximum number of changing rounds in the widening scheme and the resulting coarse layout-work bound.
\end{enumerate}

\section{Further reading}
\label{sec:ch55-reading}
Salomon treats macro expansion, conditional assembly, and source-level assembler facilities in detail~\cite{salomon-1993-assemblers}.
The Lua manual defines strings, integer behavior, iteration, and the C API needed to embed this preprocessing phase~\cite{ierusalimschy-2026-lua}.
GNU linker documentation illustrates the related layout and relaxation responsibilities that can remain after assembly~\cite{gnu-2026-ld}.
